%% PACKAGES %%

\usepackage{makeidx}
\usepackage[english,italian]{babel}
\usepackage{floatflt}
\usepackage{titlesec}
\usepackage{listings}
\usepackage{amsmath}
\usepackage{amsthm}
\usepackage{amsfonts}
\usepackage{hyperref}
\usepackage{graphics}
\usepackage{wrapfig}
\usepackage[utf8]{inputenc}
\usepackage{subfiles}
\usepackage[version=3]{mhchem}
\usepackage{pgfplots}
\usepackage{amssymb}
\usepackage{textcomp} 
\usepackage{enumitem}
\usepackage{blindtext}
\usepackage{xcolor}
\usepackage{mathtools}
\usepackage[thinc]{esdiff}
\usepackage[T1]{fontenc}
\usepackage[ddmmyyyy]{datetime}

\usepackage{caption}
\usepackage{chngcntr}
\counterwithin{figure}{section}
\numberwithin{equation}{section}
\counterwithin{table}{section}

\usepackage{circuitikz}

%% TEST PACKAGE %%

\usepackage{fancyhdr}
\pagestyle{fancy}
\fancyhf{}

\setlength{\headwidth}{16cm}
\setlength{\headheight}{14pt}

% PAGE NUMERATION

\fancyhead[LE]{\thepage}
\fancyhead[RO]{\thepage}

\fancyhead[CE]{\nouppercase{\leftmark}}
\fancyhead[CO]{\nouppercase{\rightmark}}

\renewcommand{\sectionmark}[1]{\markboth{\thesection\ -\ #1}{}}
\renewcommand{\subsectionmark}[1]{\markright{\thesubsection\ -\ #1}}

\renewcommand{\headrulewidth}{0.4pt}

\fancypagestyle{plain}{
    \fancyhf{}
    \fancyhead[LE]{\thepage\quad\leftmark}
    \fancyhead[RO]{\rightmark\quad\thepage}
    \renewcommand{\headrulewidth}{0pt}
}

%% CUSTOM EQUATION SPACE %%

\AtBeginDocument{
\setlength{\abovedisplayskip}{15pt}
\setlength{\belowdisplayskip}{15pt}
\setlength{\abovedisplayshortskip}{15pt}
\setlength{\belowdisplayshortskip}{15pt}
}


%% CUSTOM CODE STYLES DISPLAY MATLAB %%

\definecolor{codegreen}{rgb}{0,0.6,0}
\definecolor{codegray}{rgb}{0.5,0.5,0.5}
\definecolor{codepurple}{rgb}{0.58,0,0.82}
\definecolor{codepurple}{rgb}{1,1,1}
\definecolor{backcolour}{rgb}{1,1,1} % BACKCOLOUR WHITE

\lstdefinestyle{matstyle}{
    backgroundcolor=\color{backcolour},   
    commentstyle=\color{codegreen},
    keywordstyle=\color{magenta},
    numberstyle=\tiny\color{codegray},
    stringstyle=\color{codepurple},
    basicstyle=\ttfamily\footnotesize,
    upquote=true,
    breakatwhitespace=false,         
    breaklines=true,                 
    captionpos=b,                    
    keepspaces=true,                 
    numbers=left,                    
    numbersep=5pt,                  
    showspaces=false,                
    showstringspaces=false,
    showtabs=false,                  
    tabsize=2
}

\lstset{style=matstyle} % STYLE FRO MATLAB CODE


%% CUSTOM CODE STYLES DISPLAY C %%

\definecolor{codebg}{RGB}{245,245,245}
\definecolor{codegreen}{rgb}{0,0.55,0}
\definecolor{codegray}{rgb}{0.5,0.5,0.5}
\definecolor{codepurple}{rgb}{0.58,0,0.82}

\lstdefinestyle{cstyle}{
    backgroundcolor   = \color{codebg},
    commentstyle      = \color{codegreen}\itshape,
    keywordstyle      = \color{blue}\bfseries,
    numberstyle       = \tiny\color{codegray},
    stringstyle       = \color{codepurple},
    basicstyle        = \ttfamily\small,
    breakatwhitespace = false,
    breaklines        = true,
    captionpos        = b,
    keepspaces        = true,
    numbers           = left,
    numbersep         = 8pt,
    showspaces        = false,
    showstringspaces  = false,
    showtabs          = false,
    tabsize           = 4,
    language          = C,
    frame             = single,
    framerule         = 0.4pt,
    rulecolor         = \color{gray},
}
\lstset{style=cstyle}



%% THEOREM STYLE SECTION %%

\newtheorem*{theorem*}{Teorema}
\newtheorem*{lemma}{Lemma}
\newtheorem*{corollary}{Corollario}
\newtheorem*{prop*}{Proprietà}
\newtheorem*{definition*}{Definizione}


%% SECTION FOR CUSTOM PATHS %%

\graphicspath{ {./images/} }

%% LAYOUT SECTION %%

\usepackage[a4paper,
    textwidth=15.5cm,
    textheight=23.0cm,
    top=2.75cm,
    inner=2.5cm,
    outer=2.5cm
]{geometry}

%% SECTION FOR CUSTOM COMMAND %%

\newcommand{\unit}[1]{\ensuremath{\, \mathrm{#1}}}
\renewcommand{\dateseparator}{-}

%% FILL SECTION FOR TITLE

\newcommand{\documentlanguage}{english}
\ifthenelse{\equal{\documentlanguage}{italian}}
    {\AtBeginDocument{\selectlanguage{italian}}}
    {\AtBeginDocument{\selectlanguage{english}}}
    
\def\AnnoS{20XX}
\def\AnnoF{20YY}
\def\StudentID{s123456}
\def\Sem{I}

%% LANGUAGE COMMAND DEPENDENT %%

\newcommand{\CourseTitle}{\iflanguage{english}{Course Title}{Titolo Corso}}
\newcommand{\CourseCode}{\iflanguage{english}{Course Code}{Codice Corso}}
\newcommand{\Professor}{\iflanguage{english}{Prof. Name Surname}{Prof. Nome Cognome}}
\newcommand{\StudentName}{\iflanguage{english}{Name}{Nome}}
\newcommand{\StudentSurname}{\iflanguage{english}{Surname}{Cognome}}
\newcommand{\Year}{\iflanguage{english}{Academic Year \AnnoS \;- \AnnoF}{Anno Accademico \AnnoS \;- \AnnoF}}
\newcommand{\Semester}{\iflanguage{english}{Semester \Sem}{Semestre \Sem}}
\newcommand{\Version}{\iflanguage{english}{Version: \version \;- \today}{Versione: \version \;- \today}}

%% TITLE INFORMATION AND VERSIONING %%

\title{\textbf{\CourseTitle \\ \CourseCode \\ \Professor}}
\author{\StudentName \; \StudentSurname \;- \StudentID}
\date{\Year \:  \Semester \\ \vspace{10pt} \small {\Version}}


% The Course title is indicated as the one listed on the PoliTO online pages
% The Course code si also listed on the university page
% The professor name will list all of the professor (multiple lines if more than one)
% The Student Name will show the student name
% The Student Id will show the ID of the student
% Version will show the current file version
% The date will be set as the date the underlying OS is using
% The Semester or year will indicate the period of the course























